\uuid{nHUq}
\exo7id{5861}
\titre{exo7 5861}
\auteur{rouget}
\organisation{exo7}
\datecreate{2010-10-16}
\isIndication{false}
\isCorrection{true}
\chapitre{Topologie}
\sousChapitre{Application linéaire continue, norme matricielle}
\cpge{mpsi-mp2i.calcul-matriciel-systemes.a, pcsi.calcul-matriciel-systemes-lineaires.a, ptsi.calcul-matriciel-systemes-lineaires.a, mpsi-mp2i.calcul-matriciel-systemes.e, pcsi.determinants.c, ptsi.determinants.a, bcpst1.matrices, mp-mpi.topologie-espaces-normes.a, pc.espaces-vectoriels-normes.a, psi.espaces-vectoriels-normes.a, bcpst2.applications-lineaires-matrices.a}
\module{Analyse}
\niveau{L2}
\difficulte{}

\contenu{
\question{
Existe-t-il une norme $N$ sur $\mathcal{M}_n(\Rr)$ ($n\geqslant2$) telle que $\forall(A,B)\in(\mathcal{M}_n(\Rr))^2$, $N(AB)=N(A)N(B)$.
}
\reponse{
Non, car si $A=E_{1,1}\neq0$ et $B=E_{2,2}\neq0$ alors $AB=0$ puis $N(AB)<N(A)N(B)$.
}
}
